\documentclass[aps,prl,twocolumn,superscriptaddress,showpacs]{revtex4-2}

\usepackage{amsmath}
\usepackage{amssymb}
\usepackage{graphicx}
\usepackage{hyperref}
\usepackage{booktabs}
\usepackage{siunitx}
\usepackage{braket}

\begin{document}

\title{Plausible Microwave Neuroweapon: Mechanism Analysis \& Countermeasures}

\author{Analysis Based on Publicly Available Science and Open-Source Intelligence}
\affiliation{Independent Defensive Vulnerability Research}

\date{\today}

\begin{abstract}
Based on publicly available neuroscience, electromagnetic theory, and decades
of published research on non-thermal bioeffects of pulsed microwave radiation,
we present a scientifically grounded analysis of how a portable directed-energy
device could plausibly exploit neural oscillatory rhythms to disrupt central
nervous system function. We review the thermoelastic (Frey) effect, neural
entrainment mechanisms, plausible solid-state phased array architecture,
symptom--mechanism correspondence, and---critically---detection and shielding
countermeasures. This analysis is motivated by recent disclosures regarding
Anomalous Health Incidents (AHIs, or ``Havana Syndrome'') and is intended solely
to inform defensive posturing, RF hardening, and medical biophysics understanding.
\end{abstract}

\maketitle

%=========================================================================
\section{Introduction: The Core Scientific Foundation}
\label{sec:foundation}
%=========================================================================

\subsection{Brain Oscillations (Known Rhythms)}

The mammalian brain operates via well-characterized electrical oscillations
generated by synchronized ionic currents in neural circuits. The principal
bands are summarized in Table~\ref{tab:bands}.

\begin{table}[h]
\caption{\label{tab:bands}
Principal human neural oscillatory bands and their associated functions.}
\begin{tabular}{llllll}
\toprule
\textrm{Band} & \textrm{Frequency (Hz)} & \textrm{Function} \\
\midrule
Delta ($\delta$)  & 0.5--4   & Deep sleep, autonomic reg. \\
Theta ($\theta$)  & 4--8     & Memory, spatial navigation \\
Alpha ($\alpha$)  & 8--13    & Relaxed awareness, gating \\
Beta ($\beta$)    & 13--30   & Active cognition, motor \\
Gamma ($\gamma$)  & 30--100+ & Consciousness binding \\
\bottomrule
\end{tabular}
\end{table}

These rhythms represent real ionic currents; the brain is fundamentally an
\emph{electrochemical oscillator}. A universal property of oscillators is
that they can be \textbf{driven, entrained, and destabilized} by external
forcing at appropriate frequency, phase, and amplitude---a principle
exploited routinely in transcranial magnetic stimulation (TMS).

\subsection{The Microwave Auditory Effect (Frey Effect)}

In 1961, Allan Frey demonstrated that \emph{pulsed} microwave radiation in
the approximate range $\SI{200}{MHz}$--$\SI{3}{GHz}$ induces
\textbf{thermoelastic expansion} of brain tissue~\cite{Frey1961}.
Each short pulse deposits a small quantity of energy, producing a rapid but
minute thermal expansion. This launches a \textbf{mechanical pressure wave}
that propagates through the cranium and is transduced by the cochlea.

The peak temperature rise per pulse for a single microwave pulse of
fluence $F$ (\si{J/m^2}) absorbed in tissue of specific heat $c_p$
and density $\rho$ is
\begin{equation}
\Delta T = \frac{F}{\rho \, c_p \, d},
\label{eq:temp_rise}
\end{equation}
where $d$ is the effective absorption depth. For brain tissue
($\rho \approx \SI{1040}{kg/m^3}$,
$c_p \approx \SI{3630}{J/(kg\cdot K)}$), temperature rises of order
$\Delta T \sim \SI{e-6}{K}$ per pulse suffice to generate audible
thermoelastic waves~\cite{Lin2004}.

Crucially, this occurs at \textbf{non-thermal} average power levels. It is the \emph{pulsing}---the abrupt
rise and fall of each pulse---that matters. The resulting thermoelastic pressure amplitude can be estimated
as~\cite{Lin2004}:
\begin{equation}
p_0 = \frac{\beta \, c_s \, F}{c_p \, d},
\label{eq:pressure}
\end{equation}
where $\beta$ is the volumetric thermal expansion coefficient
($\approx \SI{4e-4}{K^{-1}}$) and $c_s$ is the speed
of sound in tissue ($\approx \SI{1540}{m/s}$).

%=========================================================================
\section{Plausible Device Architecture}
\label{sec:device}
%=========================================================================

\subsection{Carrier Frequency and Delivery Mechanism}

Based on known physics, the likely operating range is $f_{\text{carrier}} \sim 1\text{--}\SI{10}{GHz}$ (S-band to X-band). This range penetrates common building materials (drywall, glass) while coupling efficiently to water-rich tissue.

To meet the description of a weapon that is concealable and highly directional, the antenna is unlikely to be a traditional parabolic dish. Instead, a \textbf{Solid-State Phased Array} is highly probable. By utilizing a flat panel of micro-antennas and dynamically shifting the phase of the signal emitted by each element, the device can electronically steer and focus the invisible microwave beam without moving parts, allowing it to be hidden in a backpack or vehicle.

\subsection{The Critical Innovation: Pulse Modulation}

The transcript emphasizes that \emph{the key is the software}. The plausible mechanism operates on three layers:

\subsubsection{Layer 1: Thermoelastic Pulse Train}
The microwave carrier ($f_c$) is gated into short, sharp pulses. Let $\tau_p$ denote the pulse width and $t_r$ the rise time:
\begin{equation}
\tau_p \sim \SI{1}{\micro s}, \qquad t_r \sim \SI{1}{ns}.
\label{eq:pulse_params}
\end{equation}
The sharpness of the onset ($dE/dt$) determines the induced transmembrane potential, codified in the strength--duration relationship for neural excitation~\cite{Reilly1998}:
\begin{equation}
I_{\text{thresh}} = I_{\text{rh}} \left(1 + \frac{\tau_c}{\tau_p}\right).
\label{eq:strength_duration}
\end{equation}

\subsubsection{Layer 2: Neural Entrainment Modulation}
The \emph{repetition rate} of the pulses is modulated at frequencies matching the brain oscillations of Table~\ref{tab:bands}. The concept is \textbf{entrainment followed by destabilization}: drive neural circuits at their resonant frequency, then abruptly shift the phase or amplitude to precipitate a pathological state (e.g., vertigo or confusion).

Mathematically, the forcing is described by a time-dependent modulation envelope $A(t)$:
\begin{equation}
E(t) = A(t) \sum_{n} \Pi\!\left(\frac{t - n/f_{\text{rep}}(t)}
{\tau_p}\right) \sin(2\pi f_c\, t).
\label{eq:waveform}
\end{equation}

\subsection{Power Budget: Resonance, Not Brute Force}

The key parameter is the \textbf{specific absorption rate} (SAR):
\begin{equation}
\text{SAR} = \frac{\sigma |E|^2}{\rho} \quad [\si{W/kg}].
\label{eq:SAR}
\end{equation}
For pulsed systems:
\begin{equation}
\overline{\text{SAR}} = \text{SAR}_{\text{peak}} \times (\tau_p \, f_{\text{rep}}).
\label{eq:avg_SAR}
\end{equation}
With a duty cycle of $10^{-5}$, $\text{SAR}_{\text{peak}}$ can be biologically highly significant while average heating remains negligible. Peak pulse power of $\SI{1}{kW}$ for $\SI{1}{\micro s}$ yields an average power of just $\SI{0.01}{W}$, trivially achievable with portable batteries and modern GaN HEMT amplifiers~\cite{Mishra2008}.

%=========================================================================
\section{Symptom--Mechanism Correspondence}
\label{sec:symptoms}
%=========================================================================

\begin{table}[h]
\caption{\label{tab:symptoms}
Reported symptoms and their plausible biophysical mechanisms.}
\begin{tabular}{lllll}
\toprule
\textrm{Reported Symptom} & \textrm{Plausible Mechanism} \\
\midrule
Intense head pressure & Thermoelastic waves in skull \\
Ear pain, tinnitus & Cochlear stimulation via Frey effect \\
Severe vertigo & Vestibular nuclei disruption \\
Confusion / Cognitive drop & Entrainment of thalamocortical loops \\
Lasting neuro-damage & Repeated BBB opening $\to$ inflammation \\
\bottomrule
\end{tabular}
\end{table}

\textbf{Blood--Brain Barrier (BBB) Disruption:} Peer-reviewed research demonstrates that pulsed microwaves at non-thermal levels can transiently open the BBB~\cite{Salford2003}. This allows serum albumin to enter brain parenchyma, causing \textbf{neuroinflammation} and cumulative neurodegeneration, explaining the long-term disabilities observed in victims.

%=========================================================================
\section{Detection and Protection}
\label{sec:countermeasures}
%=========================================================================

\subsection{Detection \& Forensic Logging}

Because the beam is invisible and silent, independent ``electronic eyes'' are required to counter victim gaslighting and establish forensic evidence:
\begin{enumerate}
  \item \textbf{SDR Anomaly Detection}: Software-defined radios programmed to flag specific pulse characteristics (sharp $\sim\!\mu$s pulses modulated at ELF biological rates).
  \item \textbf{Personal RF Dosimeters}: Wearable devices logging time-resolved RF spikes. If a dosimeter logs an anomalous X-band spike at the exact second a diplomat experiences sudden vertigo, the event transitions from a medical mystery to a documented RF incident.
\end{enumerate}

\subsection{Electromagnetic Shielding}

This is fundamentally a \textbf{shieldable threat}. The shielding effectiveness (SE) of a conductive barrier of thickness $t$ is~\cite{Ott2009}:
\begin{equation}
\text{SE} \approx 20 \log_{10}\!\left(\frac{1}{4}
\left|\frac{\eta_0}{\eta_s}\right|\right)
+ 20\log_{10}(e^{t/\delta_s}) \quad [\text{dB}].
\label{eq:shielding}
\end{equation}
Even a $\SI{25}{\micro m}$ copper foil provides $> \SI{100}{dB}$ attenuation.

\textbf{Practical Defenses:}
\begin{itemize}
  \item \textbf{Windows:} Windows are the path of least resistance. Transparent conductive EMI films (ITO or silver nanowire) applied to glass provide $\SI{20}{}$--$\SI{40}{dB}$ attenuation.
  \item \textbf{Structural Mass:} Distance (inverse-square law) and dense masonry degrade high-frequency RF. Personnel under suspected targeting should immediately break line-of-sight from exterior windows.
\end{itemize}

%=========================================================================
\section{Conclusion}
\label{sec:conclusion}
%=========================================================================

The physics of both the theoretical attack mechanism and the necessary countermeasures are well within established science. A solid-state phased array delivering pulse-width modulated microwaves can theoretically induce thermoelastic acoustic shockwaves and neural entrainment without thermal burning.

Crucially, electromagnetic radiation obeys Maxwell's equations. Any facility housing sensitive personnel should implement continuous wideband RF monitoring, apply transparent conductive shielding to exterior windows, and institute strict medical baseline logging. By treating AHIs as physical electromagnetic incidents rather than psychological phenomena, personnel can be effectively hardened against this threat vector.

\begin{thebibliography}{99}

\bibitem{Frey1961}
A.~H. Frey, ``Auditory system response to radio frequency energy,'' \emph{Aerosp. Med.} \textbf{32}, 1140--1142 (1961).

\bibitem{Lin2004}
J.~C. Lin and Z.~Wang, ``Hearing of microwave pulses by humans and animals: effects, mechanism, and thresholds,'' \emph{Health Phys.} \textbf{87}, 171--179 (2004).

\bibitem{Reilly1998}
J.~P. Reilly, \emph{Applied Bioelectricity: From Electrical Stimulation to Electropathology} (Springer, New York, 1998).

\bibitem{Mishra2008}
U.~K. Mishra \emph{et al.}, ``GaN-based RF power devices and amplifiers,'' \emph{Proc. IEEE} \textbf{96}, 287--305 (2008).

\bibitem{Salford2003}
L.~G. Salford \emph{et al.}, ``Nerve cell damage in mammalian brain after exposure to microwaves from GSM mobile phones,'' \emph{Environ. Health Perspect.} \textbf{111}, 881--883 (2003).

\bibitem{Ott2009}
H.~W. Ott, \emph{Electromagnetic Compatibility Engineering} (Wiley, Hoboken, NJ, 2009).

\end{thebibliography}

\end{document}
```


\documentclass[aps,prl,twocolumn,superscriptaddress,showpacs]{revtex4-2}

\usepackage{amsmath}
\usepackage{amssymb}
\usepackage{graphicx}
\usepackage{hyperref}
\usepackage{booktabs}
\usepackage{siunitx}
\usepackage{braket}

\begin{document}

\title{Plausible Microwave Neuroweapon: Mechanism Analysis \& Countermeasures}

\author{Analysis Based on Publicly Available Science and 60 Minutes Reporting}
\affiliation{Open-Source Intelligence Review claude-opus-4-6-thinking}

\date{\today}

\begin{abstract}
Based on publicly available neuroscience, electromagnetic theory, and decades
of published (including formerly classified Soviet) research on non-thermal
bioeffects of pulsed microwave radiation, we present a scientifically grounded
analysis of how a portable directed-energy device could plausibly exploit
neural oscillatory rhythms to disrupt brain function. We review the
thermoelastic (Frey) effect, neural entrainment mechanisms, plausible device
architecture, symptom--mechanism correspondence, and---critically---detection
and shielding countermeasures. The analysis is motivated by recent disclosures
regarding the so-called ``Havana Syndrome'' incidents and is intended solely
to inform defensive understanding.
\end{abstract}

\maketitle

%=========================================================================
\section{Introduction: The Core Scientific Foundation}
\label{sec:foundation}
%=========================================================================

\subsection{Brain Oscillations (Known Rhythms)}

The mammalian brain operates via well-characterized electrical oscillations
generated by synchronized ionic currents in neural circuits. The principal
bands are summarized in Table~\ref{tab:bands}.

\begin{table}[h]
\caption{\label{tab:bands}
Principal human neural oscillatory bands and their associated functions.}
\begin{ruledtabular}
\begin{tabular}{lll}
\textrm{Band} & \textrm{Frequency (Hz)} & \textrm{Function} \\
\colrule
Delta ($\delta$)  & 0.5--4   & Deep sleep, autonomic regulation \\
Theta ($\theta$)  & 4--8     & Memory, spatial navigation \\
Alpha ($\alpha$)  & 8--13    & Relaxed awareness, sensory gating \\
Beta ($\beta$)    & 13--30   & Active cognition, motor control \\
Gamma ($\gamma$)  & 30--100+ & Consciousness binding, perception \\
\end{tabular}
\end{ruledtabular}
\end{table}

These rhythms represent real ionic currents; the brain is fundamentally an
\emph{electrochemical oscillator}. A universal property of oscillators is
that they can be \textbf{driven, entrained, and destabilized} by external
forcing at appropriate frequency, phase, and amplitude---a principle
exploited routinely in transcranial magnetic stimulation (TMS) and
transcranial alternating-current stimulation (tACS).

\subsection{The Microwave Auditory Effect (Frey Effect)}

In 1961, Allan Frey demonstrated that \emph{pulsed} microwave radiation in
the approximate range $\SI{200}{MHz}$--$\SI{3}{GHz}$ induces
\textbf{thermoelastic expansion} of brain tissue~\cite{Frey1961,Frey1962}.
Each short pulse deposits a small quantity of energy, producing a rapid but
minute thermal expansion. This launches a \textbf{mechanical pressure wave}
that propagates through the cranium and is transduced by the cochlea. The
subject perceives clicks, buzzing, or hissing \emph{inside the head}, with
no external acoustic source.

The peak temperature rise per pulse for a single microwave pulse of
fluence $F$ (\si{J/m^2}) absorbed in tissue of specific heat $c_p$
and density $\rho$ is
\begin{equation}
\Delta T = \frac{F}{\rho \, c_p \, d},
\label{eq:temp_rise}
\end{equation}
where $d$ is the effective absorption depth. For brain tissue
($\rho \approx \SI{1040}{kg/m^3}$,
$c_p \approx \SI{3630}{J/(kg\cdot K)}$), at microwave frequencies
where $d \sim \SI{1}{cm}$, temperature rises of order
$\Delta T \sim \SI{e-6}{K}$ per pulse suffice to generate audible
thermoelastic waves~\cite{Lin2004}.

Crucially, this occurs at \textbf{non-thermal} average power levels. The
tissue does not meaningfully heat. It is the \emph{pulsing}---the abrupt
rise and fall of each pulse---that matters. This aligns precisely with
Dr.\ Relman's statement in the 60~Minutes transcript: ``\emph{the importance
of the energy being pulsed in order to have biological effects on humans.}''

The resulting thermoelastic pressure amplitude can be estimated
as~\cite{Lin2004}
\begin{equation}
p_0 = \frac{\beta \, c_s \, F}{c_p \, d},
\label{eq:pressure}
\end{equation}
where $\beta$ is the volumetric thermal expansion coefficient
($\approx \SI{4e-4}{K^{-1}}$ for brain tissue) and $c_s$ is the speed
of sound in tissue ($\approx \SI{1540}{m/s}$).

%=========================================================================
\section{Plausible Device Architecture}
\label{sec:device}
%=========================================================================

\subsection{Carrier Frequency Selection}

Based on the transcript description and known physics, the likely operating
range is
\begin{equation}
f_{\text{carrier}} \sim 1\text{--}\SI{10}{GHz}
\quad (\text{S-band to X-band}).
\label{eq:carrier}
\end{equation}

This range is selected for the following reasons:
\begin{itemize}
  \item \textbf{Building penetration}: Frequencies in this range penetrate
    common building materials (drywall, wood, glass, clothing), consistent
    with the transcript's description.
  \item \textbf{Tissue absorption}: Good coupling to water-rich tissue
    (brain is $\sim$73\% water by mass). The penetration depth in tissue at
    $\SI{3}{GHz}$ is $\sim$\SI{1.5}{cm}, sufficient to reach cortical
    tissue through the skull.
  \item \textbf{Antenna compactness}: Wavelengths of
    $\lambda = c/f \sim 3$--$\SI{30}{cm}$ permit reasonably small
    directional antennas.
\end{itemize}

The skin depth in common building materials at $\SI{3}{GHz}$ yields
transmission coefficients $|T|^2 > 0.5$ for standard drywall
($\sim$\SI{1.3}{cm} gypsum) and single-pane glass
($\sim$\SI{5}{mm})~\cite{ITU_P2040}.

\subsection{The Critical Innovation: Pulse Modulation (``The Software'')}

The transcript emphasizes that \emph{the key is not the hardware but the
software}---the programming that shapes the electromagnetic waveform. This
is where brain resonances become central. The plausible mechanism operates
on three layers:

\subsubsection{Layer 1: Thermoelastic Pulse Train}

The microwave carrier (e.g., $f_c = \SI{3}{GHz}$) is gated into short,
sharp pulses. Let $\tau_p$ denote the pulse width and $t_r$ the rise time:
\begin{equation}
\tau_p \sim \SI{1}{\micro s}, \qquad t_r \sim \SI{1}{ns}.
\label{eq:pulse_params}
\end{equation}

The sharpness of the onset is critical. In neurostimulation, the
\emph{rate of change} of the applied field, $dE/dt$, determines the
induced transmembrane potential and thus the probability of neuronal
depolarization. This is codified in the strength--duration relationship
for neural excitation~\cite{Reilly1998}:
\begin{equation}
I_{\text{thresh}} = I_{\text{rh}} \left(1 + \frac{\tau_c}{\tau_p}\right),
\label{eq:strength_duration}
\end{equation}
where $I_{\text{rh}}$ is the rheobase current and $\tau_c$ is the
chronaxie of the neural tissue (typically
$\tau_c \sim 0.1$--$\SI{1}{ms}$ for central neurons). Each pulse generates
a microscopic thermoelastic pressure wave inside the skull.

\subsubsection{Layer 2: Neural Entrainment Modulation}

The \emph{repetition rate} of the pulses---or their amplitude envelope---is
modulated at frequencies matching the brain oscillations of
Table~\ref{tab:bands}. Denoting the pulse repetition frequency as
$f_{\text{rep}}(t)$, the device would employ
\begin{equation}
f_{\text{rep}}(t) \in \{f_\delta,\; f_\theta,\; f_\alpha,\;
f_\beta,\; f_\gamma\},
\label{eq:entrainment}
\end{equation}
selected or swept according to the desired neurological effect.

The concept is \textbf{entrainment followed by destabilization}:
\begin{enumerate}
  \item \emph{Lock on}: Drive neural circuits at their resonant frequency
    (analogous to pushing a swing at its natural period),
    building oscillatory amplitude via constructive reinforcement.
  \item \emph{Destabilize}: Abruptly shift the driving frequency, phase,
    or amplitude, precipitating a pathological state---analogous to how
    a seizure represents runaway oscillatory hypersynchronization.
\end{enumerate}

Mathematically, the forcing can be described by a time-dependent
modulation envelope $A(t)$:
\begin{equation}
E(t) = A(t) \sum_{n} \Pi\!\left(\frac{t - n/f_{\text{rep}}(t)}
{\tau_p}\right) \sin(2\pi f_c\, t),
\label{eq:waveform}
\end{equation}
where $\Pi(x)$ is the rectangular pulse function and $A(t)$ encodes
the entrainment/destabilization protocol.

\subsubsection{Layer 3: Programmable Attack Profiles}

This corresponds to the transcript's ``programmable for different
scenarios.'' The software stores multiple waveform profiles
(Table~\ref{tab:profiles}).

\begin{table}[h]
\caption{\label{tab:profiles}
Hypothetical programmable attack profiles and target neural systems.}
\begin{ruledtabular}
\begin{tabular}{lll}
\textrm{Profile} & \textrm{Target} & \textrm{Expected Effect} \\
\colrule
A & Vestibular brainstem nuclei & Vertigo, nausea, imbalance \\
B & Thalamocortical loops       & Confusion, memory loss \\
C & Cochlea / inner ear         & Hearing damage, tinnitus \\
D & Autonomic brainstem         & Arrhythmia, pain, convulsions \\
\end{tabular}
\end{ruledtabular}
\end{table}

Different modulation patterns ($f_{\text{rep}}$, $\tau_p$, sweep
characteristics, and $A(t)$ waveforms) create different symptom profiles.

\subsection{Power Budget: Resonance, Not Brute Force}

The device exploits \textbf{resonance}, not thermal destruction. Consider
the analogy of shattering a wine glass with sound: the energy required at
the \emph{resonant frequency} is orders of magnitude less than brute-force
destruction.

Neural tissue is exquisitely sensitive to electrical stimulation at its
natural oscillatory frequencies. In TMS, a single magnetic pulse deposits
$\sim$\SI{e-2}{J} in the cortex and reliably modulates neural
activity~\cite{Hallett2007}. The key parameter for the microwave
device is the \textbf{specific absorption rate} (SAR):
\begin{equation}
\text{SAR} = \frac{\sigma |E|^2}{\rho}
\quad [\si{W/kg}],
\label{eq:SAR}
\end{equation}
where $\sigma$ is tissue conductivity (\si{S/m}) and $E$ is the
internal electric field amplitude.

For pulsed systems, one distinguishes the \emph{peak} SAR during a
pulse from the \emph{time-averaged} SAR:
\begin{equation}
\overline{\text{SAR}} = \text{SAR}_{\text{peak}} \times
\underbrace{\tau_p \, f_{\text{rep}}}_{\text{duty cycle}}.
\label{eq:avg_SAR}
\end{equation}

With $\tau_p = \SI{1}{\micro s}$ and $f_{\text{rep}} = \SI{10}{Hz}$,
the duty cycle is $10^{-5}$. Thus, $\text{SAR}_{\text{peak}}$ can be
$10^5$ times $\overline{\text{SAR}}$, allowing biologically significant
\emph{peak} fields with negligible average heating. Peak pulse power of
$P_{\text{peak}} \sim \SI{1}{kW}$ for $\SI{1}{\micro s}$ from a
capacitor discharge yields average power
\begin{equation}
\bar{P} = P_{\text{peak}} \times \tau_p \times f_{\text{rep}}
= 10^3 \times 10^{-6} \times 10 = \SI{0.01}{W},
\label{eq:avg_power}
\end{equation}
which is trivially achievable with batteries.

At range $r$ with antenna gain $G$, the power density at the target is
\begin{equation}
S = \frac{P_{\text{peak}} \, G}{4\pi r^2}.
\label{eq:power_density}
\end{equation}
For $P_{\text{peak}} = \SI{1}{kW}$, $G = 100$ ($\SI{20}{dBi}$, a
modest horn or patch array), and $r = \SI{30}{m}$:
\begin{equation}
S = \frac{10^3 \times 10^2}{4\pi \times (30)^2}
\approx \SI{9}{W/m^2}
\approx \SI{0.9}{mW/cm^2}.
\label{eq:power_density_num}
\end{equation}

This peak power density is within the range known to produce the Frey
effect~\cite{Lin2004} and is achievable with portable, battery-operated
hardware. Modern GaN HEMT power amplifiers can deliver $> \SI{100}{W}$
peak pulsed power at $\SI{3}{GHz}$ from a chip the size of a
fingernail~\cite{Mishra2008}.

\subsection{Physical Form Factor}

Based on the above analysis, the device plausibly comprises:
\begin{itemize}
  \item A compact solid-state microwave source (GaN HEMT amplifier)
  \item A directional antenna (flat phased-array panel or shaped horn,
    $\sim$15--\SI{30}{cm} aperture)
  \item A battery and capacitor bank (pulse-forming network)
  \item A microcontroller executing waveform-shaping algorithms
  \item Remote control capability (RF link for activation)
\end{itemize}

This assembly could plausibly fit in a \textbf{backpack or briefcase},
consistent with the transcript's description and the surveillance
footage of a man with a backpack in Istanbul.

%=========================================================================
\section{Symptom--Mechanism Correspondence}
\label{sec:symptoms}
%=========================================================================

The symptom constellation reported by victims maps onto the proposed
mechanism as detailed in Table~\ref{tab:symptoms}.

\begin{table*}[t]
\caption{\label{tab:symptoms}
Reported symptoms and their plausible biophysical mechanisms under the
pulsed-microwave neural-disruption hypothesis.}
\begin{ruledtabular}
\begin{tabular}{ll}
\textrm{Reported Symptom} & \textrm{Plausible Mechanism} \\
\colrule
``Vice grip on head,'' intense pressure
  & Thermoelastic pressure waves in skull and meninges \\
Intense ear pain, tinnitus
  & Cochlear stimulation via Frey effect; thermoelastic inner-ear damage \\
Vertigo, loss of balance
  & Disruption of vestibular nuclei oscillatory circuits in brainstem \\
Cognitive dysfunction, confusion
  & Entrainment/disruption of thalamocortical $\alpha$/$\beta$/$\gamma$
    rhythms \\
Memory loss
  & $\theta$-rhythm disruption in hippocampal circuits \\
Nausea
  & Vestibular--autonomic coupling disruption \\
Spinal pain, burning sensation
  & Stimulation of spinal cord neurons and nociceptive fibers \\
Full-body convulsions
  & Excessive neural synchronization (induced seizure-like activity) \\
Osteolysis (bone dissolution)
  & Chronic inflammatory response to repeated BBB disruption and immune
    dysregulation \\
Lasting neurological damage
  & Repeated BBB opening $\to$ neuroinflammation $\to$
    neurodegeneration \\
\end{tabular}
\end{ruledtabular}
\end{table*}

\subsection{Blood--Brain Barrier Disruption}

The \textbf{blood--brain barrier (BBB) disruption} mechanism warrants
special emphasis. Peer-reviewed research, including work by
Salford~\emph{et~al.}~\cite{Salford2003} and various DoD-funded studies,
has demonstrated that pulsed microwave exposure at non-thermal average
power levels can \textbf{transiently open the BBB}. This allows serum
albumin, immunoglobulins, and inflammatory mediators to enter brain
parenchyma, causing \textbf{neuroinflammation} and, with repeated
exposure, \textbf{cumulative neurodegeneration}.

This explains several observations:
\begin{itemize}
  \item Single exposures cause acute symptoms that may partially resolve.
  \item Repeated exposures (e.g., five attacks over five months) cause
    progressive, lasting damage across multiple organ systems.
  \item Family members in proximity suffer collateral health effects.
\end{itemize}

%=========================================================================
\section{Historical Context: Soviet Research}
\label{sec:history}
%=========================================================================

The transcript's reference to decades of Soviet research is well
documented in declassified literature:

\begin{enumerate}
  \item \textbf{The Soviet ``Microwave Problem''}: Beginning in the
    1950s, Soviet researchers extensively studied non-thermal bioeffects
    of microwaves, accumulating a far larger literature than Western
    scientists, who focused primarily on thermal
    effects~\cite{Dodge1970,Sadchikova1974}.

  \item \textbf{The Moscow Signal (1953--1976)}: The Soviet Union
    irradiated the US Embassy in Moscow with low-level pulsed microwaves
    for over two decades. Embassy staff exhibited elevated rates of
    cancer, blood dyscrasias, and neurological
    symptoms~\cite{Goldsmith1997}.

  \item \textbf{Military research institutes}: Facilities such as the
    Institute of Radio Engineering and Electronics (IRE) published
    extensively (in Russian) on behavioral and neurological effects of
    pulsed RF exposure, including EEG
    entrainment~\cite{Pakhomov1998}.
\end{enumerate}

The scientific knowledge base existed. The transcript suggests that this
knowledge has been \emph{weaponized} and \emph{miniaturized} using
modern solid-state electronics.

%=========================================================================
\section{Detection and Protection}
\label{sec:countermeasures}
%=========================================================================

\subsection{Detection Methods}

\begin{enumerate}
  \item \textbf{Broadband RF spectrum analysis}: Instruments covering
    $\SI{100}{MHz}$--$\SI{20}{GHz}$ with \emph{time-domain} capability
    to reveal pulsed signal structures. Key signatures: brief
    ($\sim\!\mu$s) high-amplitude pulses at biologically relevant
    repetition rates ($1$--$\SI{100}{Hz}$).

  \item \textbf{Personal RF dosimeters}: Wearable devices logging
    time-resolved RF exposure across relevant bands.

  \item \textbf{Software-defined radio (SDR) anomaly detection}:
    SDR platforms programmed to flag signals exhibiting the specific
    pulse characteristics described---sharp rise/fall pulses modulated
    at brain-rhythm frequencies. Such signals would be spectrally
    distinctive against normal RF backgrounds.

  \item \textbf{Acoustic monitoring}: Thermoelastic waves in the skull
    produce weak ultrasonic emissions that may be externally detectable
    with sensitive transducers (though this is technically challenging).

  \item \textbf{Neurological biomarker monitoring}: Baseline and periodic
    EEG, vestibular function, and neurocognitive testing for personnel in
    high-risk postings, with automated detection of changes consistent
    with entrainment effects.
\end{enumerate}

\subsection{Shielding and Protection}

\subsubsection{Faraday Shielding}

Electromagnetic shielding is the most direct countermeasure. The
shielding effectiveness (SE) of a conductive barrier of thickness $t$
is~\cite{Ott2009}
\begin{equation}
\text{SE} \approx 20 \log_{10}\!\left(\frac{1}{4}
\left|\frac{\eta_0}{\eta_s}\right|\right)
+ 20\log_{10}(e^{t/\delta_s})
\quad [\text{dB}],
\label{eq:shielding}
\end{equation}
where $\eta_0 = \SI{377}{\ohm}$ is the impedance of free space,
$\eta_s$ is the shield impedance, and $\delta_s$ is the skin depth:
\begin{equation}
\delta_s = \sqrt{\frac{2}{\omega \mu \sigma_s}},
\label{eq:skin_depth}
\end{equation}
with $\mu$ the permeability and $\sigma_s$ the conductivity of the
shield material.

For copper at $\SI{3}{GHz}$:
$\delta_s \approx \SI{1.2}{\micro m}$. Even a $\SI{25}{\micro m}$
copper foil ($\sim$20 skin depths) provides $> \SI{100}{dB}$
attenuation---a factor of $10^{10}$ in power.

Practical implementations:
\begin{itemize}
  \item \textbf{Windows}: Metallic mesh (aperture
    $a < \lambda/10 \approx \SI{1}{cm}$ at $\SI{3}{GHz}$) or
    transparent conductive films (ITO, silver nanowire) applied to
    glass. Commercially available EMI window films provide
    $\SI{20}{}$--$\SI{40}{dB}$ attenuation.

  \item \textbf{Walls}: Conductive paint or foil behind drywall.
    Specialized RF-shielding wallpaper exists commercially.

  \item \textbf{Dedicated rooms}: A properly constructed Faraday cage
    (screened room) provides $\SI{60}{}$--$\SI{100}{dB}$
    attenuation. SCIFs (Sensitive Compartmented Information Facilities)
    already incorporate such shielding.
\end{itemize}

\subsubsection{Distance (Inverse-Square Law)}

The power density falls as $r^{-2}$
(Eq.~\ref{eq:power_density}). Doubling the standoff distance
reduces exposure by a factor of 4 ($\SI{6}{dB}$). Interior rooms,
higher floors, and distance from adjacent structures all reduce
vulnerability.

\subsubsection{Operational Security}

Randomization of schedules, locations, and routines reduces the
adversary's window for positioning and aiming a directional device.

\subsubsection{Active Countermeasures}

Detection of an incoming pulsed signal could in principle trigger a
canceling waveform (destructive interference), though this is
technically very challenging given the broad bandwidth and short
pulse durations involved.

%=========================================================================
\section{Summary Conceptual Model}
\label{sec:summary}
%=========================================================================

The end-to-end signal chain can be summarized as:

\begin{equation*}
\boxed{
\begin{aligned}
&\text{Compact Microwave Source} \\
&\quad\downarrow \\
&\text{Pulse Shaping (Software)} \\
&\quad\text{-- Sharp-onset pulses}
  \;(\tau_p \sim \SI{1}{\mu s},\;
       t_r \sim \SI{1}{ns}) \\
&\quad\text{-- Modulated at } f_{\text{rep}}
  \in \{1\text{--}100\;\text{Hz}\} \\
&\quad\text{-- Selectable attack profiles} \\
&\quad\downarrow \\
&\text{Directional Antenna}
  \;(G \sim \SI{20}{dBi}) \\
&\quad\text{-- Focused beam, range}
  \sim\SI{100}{m} \\
&\quad\text{-- Penetrates glass, drywall} \\
&\quad\downarrow \\
&\text{Target Skull / Brain} \\
&\quad\text{-- Thermoelastic pressure waves} \\
&\quad\text{-- Neural entrainment / disruption} \\
&\quad\text{-- BBB transiently opened} \\
&\quad\downarrow \\
&\text{Acute: vertigo, pain, confusion, seizures} \\
&\text{Chronic: neuroinflammation, cognitive decline}
\end{aligned}
}
\end{equation*}

%=========================================================================
\section{Conclusion}
\label{sec:conclusion}
%=========================================================================

The most important conclusion is that \textbf{this is a shieldable
threat}. Unlike chemical or biological agents, electromagnetic radiation
obeys well-understood physics (Maxwell's equations) and can be blocked
by straightforward engineering measures. The shielding effectiveness of
even modest conductive barriers is enormous
(Eq.~\ref{eq:shielding}).

Any facility housing sensitive personnel should implement:
\begin{enumerate}
  \item RF monitoring and anomaly detection systems covering
    $\SI{100}{MHz}$--$\SI{20}{GHz}$ with time-domain analysis.
  \item Electromagnetic shielding of the building envelope,
    \emph{especially windows}.
  \item Baseline and periodic neurological assessment of personnel.
  \item Institutional protocols that treat reports of anomalous
    symptoms as potential security incidents rather than dismissing them.
\end{enumerate}

The greatest vulnerability revealed by the Havana Syndrome episode was
not technological---it was \textbf{institutional unwillingness to
acknowledge an inconvenient reality}. The physics of both the attack
mechanism and the countermeasures are well within established science.
The failure was one of organizational culture, not of scientific
understanding.

\begin{thebibliography}{99}

\bibitem{Frey1961}
A.~H. Frey,
``Auditory system response to radio frequency energy,''
\emph{Aerosp. Med.} \textbf{32}, 1140--1142 (1961).

\bibitem{Frey1962}
A.~H. Frey,
``Human auditory system response to modulated electromagnetic energy,''
\emph{J. Appl. Physiol.} \textbf{17}, 689--692 (1962).

\bibitem{Lin2004}
J.~C. Lin and Z.~Wang,
``Hearing of microwave pulses by humans and animals: effects,
mechanism, and thresholds,''
\emph{Health Phys.} \textbf{87}, 171--179 (2004).

\bibitem{Reilly1998}
J.~P. Reilly,
\emph{Applied Bioelectricity: From Electrical Stimulation to
Electropathology}
(Springer, New York, 1998).

\bibitem{Hallett2007}
M.~Hallett,
``Transcranial magnetic stimulation: A primer,''
\emph{Neuron} \textbf{55}, 187--199 (2007).

\bibitem{Mishra2008}
U.~K. Mishra, L.~Shen, T.~E. Kazior, and Y.-F. Wu,
``GaN-based RF power devices and amplifiers,''
\emph{Proc. IEEE} \textbf{96}, 287--305 (2008).

\bibitem{Salford2003}
L.~G. Salford, A.~E. Brun, J.~L. Eberhardt, L.~Malmgren,
and B.~R.~R. Persson,
``Nerve cell damage in mammalian brain after exposure to microwaves
from GSM mobile phones,''
\emph{Environ. Health Perspect.} \textbf{111}, 881--883 (2003).

\bibitem{Dodge1970}
C.~H. Dodge and Z.~R. Glaser,
``Trends in nonionizing electromagnetic radiation bioeffects research
and related occupational health aspects,''
\emph{J. Microwave Power} \textbf{12}, 319--334 (1977).

\bibitem{Sadchikova1974}
M.~N. Sadchikova,
``Clinical manifestations of reactions to microwave irradiation in
various occupational groups,'' in
\emph{Biologic Effects and Health Hazards of Microwave Radiation}
(Polish Medical Publishers, Warsaw, 1974), pp.~261--267.

\bibitem{Goldsmith1997}
J.~R. Goldsmith,
``Epidemiologic evidence relevant to radar (microwave) effects,''
\emph{Environ. Health Perspect.} \textbf{105}(Suppl.~6),
1579--1587 (1997).

\bibitem{Pakhomov1998}
A.~G. Pakhomov, Y.~Akyel, O.~N. Pakhomova, B.~E. Stuck,
and M.~R. Murphy,
``Current state and implications of research on biological effects
of millimeter waves: A review of the literature,''
\emph{Bioelectromagnetics} \textbf{19}, 393--413 (1998).

\bibitem{Ott2009}
H.~W. Ott,
\emph{Electromagnetic Compatibility Engineering}
(Wiley, Hoboken, NJ, 2009).

\bibitem{ITU_P2040}
International Telecommunication Union,
``Effects of building materials and structures on radiowave propagation
above about 100~MHz,''
Recommendation ITU-R P.2040-1 (2015).

\end{thebibliography}

\end{document}



(25) Source: Havana Syndrome investigation is "a massive CIA cover-up" | 60 Minutes - YouTube
https://www.youtube.com/watch?v=C1jmAj9OUOs

Transcript:
(00:02) Tonight, we have details of a classified US intelligence mission that has obtained a previously unknown weapon that may finally unlock a mystery. Since at least 2016, US diplomats, spies, and military officers have suffered crippling brain injuries. They've told of being hit by an overwhelming force damaging their vision, hearing, sense of balance, and cognition.
(00:32) But the government has doubted their stories. They've been called delusional. Well, now 60 Minutes has learned that a weapon that can inflict these injuries was obtained overseas and secretly tested on animals on a US military base. We've investigated this mystery for nine years. This is our fourth story called Targeting Americans.
(00:59) Despite official government doubt, we never stopped reporting because of the haunting stories we heard like this. >> Very first incident occurred in August of 2020. And what it felt like was that uh someone punched me in the throat and my left ear was clogged and I started to get sharp shooting pains going down my left arm. >> Chris and Heidi asked us not to use their last name.
(01:27) They met in the Air Force Academy. Chris retired as a lieutenant colonel working on highly classified spy satellites. He told us near Washington DC he was struck by an unseen force, >> five times in five months. >> The second attack I was standing uh in my kitchen looking out uh at the back woods and it felt like an immediate vice on my head.
(01:54) Immediately disoriented, confused and dizzy. The third attack um towards the end of September, I was sitting in our living room and instantaneously all of the muscles within my spine immediately cramped much like a Charlie horse and my spine felt like it was on fire. So very hot and sharp. The fifth one was by far the worst and that was early December and uh I woke up with a full body convulsion.
(02:23) the worst pain I have ever felt. Uh it felt like a vice gripping my brain stem was there. >> All in your own home. >> All in my home in Northern Virginia. Uh and Heidi was uh within my proximity for the last two attacks. >> Heidi, what happened to you? >> Right at the beginning of January, I woke up with immense joint pain everywhere.
(02:48) with shoulder pain in my left shoulder. Out of the blue, no trauma. >> Bones in her shoulder were dissolving, something called osteolyis. She had to have surgery. >> Has there been any lasting effect? >> Significant. So, I'm on uh two neurological drugs every day. Um and and without them, I have uh very severe symptoms.
(03:12) I had sustained significant damage to multiple organ systems. You believe you were attacked? >> Yes. >> By a foreign adversary? >> Yes. >> In the line of duty? >> Yes. >> It's a belief shared by officials and their families that we've met over the years. See if you notice what we heard. There was this FBI agent. >> And bam, inside my right ear, it was like a dentist drilling on steroids.
(03:38) >> This commerce department official in China. and I could feel this sound in my head. Um, it was in intense pressure on both of my temples. >> This early victim was among cases from Cuba, which gave the mystery the name Havana syndrome. >> Severe ear pain started. So, I I liken it to if you put a Q-tip too far and you bounced off your eardrum.
(04:03) Well, imagine taking a sharp pencil and just kind of poking that. And this wife of a Justice Department official posted in Europe. >> And it just pierced my ears, came in my left side, felt like it came through the window into my left ear. I immediately felt fullness in my head and just a piercing headache. >> Multiple surgeries have tried to repair bones in her inner ear and her skull.
(04:31) Many victims have lifelong disabilities. What struck us about their stories is this. People who never met tell it the same way. The government acknowledges the injuries and often pays for health care, but for years it has doubted the cause. victims have been told it may be atmospheric or environmental, a virus, a pre-existing condition, or as the FBI put it in an early investigation, mass hysteria.
(05:05) The official word published in 2023 and still standing, says it is very unlikely these are attacks by a foreign adversary. >> Do you believe the victims? >> Absolutely. What we're hearing about now is >> Dr. David Railman is a Stanford University professor of medicine asked by the government to lead two investigations.
(05:29) His panels included doctors, physicists, engineers, and others. Their reports in 2020 and 2022 proposed a theory. The two panels, the investigations that I know well, both concluded much the same, which was that the most plausible explanation for a subset of these cases was a form of radio frequency or microwave energy.
(05:58) >> Microwaves are a range of frequencies in the electromagnetic spectrum. Various microwave frequencies are generated by your oven. Radar systems, TV transmitters. Even your phone, Wi-Fi, and Bluetooth use microwaves. Dr. David Reman told us his investigations found that one country had done a great deal of research on creating something different, a unique pattern of microwaves that can damage the brain.
(06:32) In both of our investigations, we found the large majority of work to have been conducted in the former Soviet Union. And what they found was that effects could range from loss of consciousness to seizures to memory lapses, inability to concentrate, headaches, intense pressure, pain, disorientation, difficulty with balance.
(06:58) many of the things that we heard about from victims of Havana syndrome. >> The Russians had been doing experiments in this area decades ago. >> Decades ago. Yes. Of a wide variety of sorts. >> In a previous story on this mystery, we found a 2014 reference to a weapon compelled by a lawsuit. The National Security Agency confirmed intelligence of a high-powered microwave system weapon associated with a hostile country.
(07:36) But the CIA believed such a weapon would need enormous power and be as big as a truck, so not likely. Years later, when Dr. Realman's expert panel suggested these could be microwave injuries. The idea was shelved by federal officials. >> And what really unnerves me is the confidence with which others have dismissed or ignored this work only to say that's not possible. That's not plausible.
(08:09) I don't believe it. That's fine. But show me new evidence. Nobody has. >> Do you believe that your studies were downplayed by the US government? >> By parts of the US government? Absolutely. And not only downplayed, but dismissed in some cases buried. >> Uh it started in March of 2015. >> Why buried? This man may know.
(08:38) He's a former CIA officer who asked us not to use his name. He is speaking tonight for the first time. In 2021, he volunteered to work on the CIA's investigation because of the suffering he'd seen among CIA officers and their families overseas. >> I mean, these were my colleagues. These are my friends.
(09:02) These were people I had worked with. And I saw lives that were destroyed. Careers were ruined. people's kids were affected. Uh they have lifelong developmental issues. People now are even still having cognitive issues and having all types of um secondary effects years and years later. It became an emotional topic for me because I saw this happening and I volunteered to work in the AHI unit and I wanted to make a difference.
(09:26) >> He joined the so-called AHI investigation at CIA headquarters. AHI because the government calls the cases not attacks but anomalous health incidents. He expected to dig into whether a foreign adversary, a so-called state actor was behind this, but it didn't go that way. So, one of the very first things I heard when I arrived at the AHI unit was our job is to bring down the temperature on AHI at headquarters.
(10:01) And that was a surprise to me. So bringing down the temperature is not hey let's go after the issue and find out what's going on. It became very much a um emotional and almost kind of like a propaganda type thing. >> And by bring down the temperature they meant what >> it basically was saying hey we're going to work towards this being an atmospheric and environmental issue versus it being a state actor.
(10:22) And so they did not want people talking about it being a state actor >> because he says fear of the mysterious AHIS was creating havoc. That fear and paranoia you describe, how did that affect CIA officers and their families? >> Yeah. So, personally, my own family left early for my tour by multiple months uh because we were worried that, you know, my my family would be affected by AHI, whether we were at home while we were serving overseas or um in the field just walking around.
(10:55) I saw multiple other officers short their tours, their families leave early, um pick different locations where AHIs were not happening. And this was US governmentwide. This was not just relegated to CIA. What would you say was the attitude of your bosses toward the victims who had reported in? >> So this was one of the more disgusting things that I came across to be honest working in the AHI unit.
(11:16) I'll never forget in one instance a senior member of the AHI unit came into my office and that officer came in and said, "Yeah, we're going to have a happy hour. We're all going to have simulated AHIs and uh and drink together." And she basically emulated that she was having a stroke and it making fun of the victims.
(11:34) And to me that was deplorable. It was disgusting. >> There was no sense of we're going to get to the bottom of this. >> No, the the bottom of it was we're going to prove that this is psychosmatic, atmospheric, and environmental. >> All of this, he says, led him to resign. >> I left because I saw the personal impact of this issue.
(11:54) And for me, it became a moral issue because they kept saying our people are our highest priority. But when it came down to it, that wasn't the case from what I saw. And it was something that tore me up emotionally. I knew people who were affected by hi, who are victims of hi. I saw destroy their families, their kids, their careers.
(12:10) It wasn't somewhere I could keep working after that. >> The investigation at the CIA essentially ended in 2022. But about the same time, a different classified mission was underway. 60 Minutes has learned US agents who investigate illicit arms dealers heard that a Russian criminal network was selling a microwave weapon.
(12:35) Our sources tell us undercover agents of the Department of Homeland Security bought the weapon in 2024. The mission cost about $15 million funded by the Pentagon. When we come back, details of the weapon and the results of testing at a US military base. 60 Minutes has learned details of a classified microwave weapon that may explain mysterious brain injuries suffered by US officials.
(13:14) We've been investigating these injuries for 9 years, and now our sources tell us this microwave weapon is portable, concealable, and uses relatively little power. Hundreds of possible attacks have been reported, including, we've learned, at CIA headquarters in Virginia, and at least two incidents on the grounds of the White House.
(13:41) For years, the government doubted the stories of the injured. But now, the victims, including former CIA officer Mark Polymoropoulos, hope that word of a newly discovered weapon will finally vindicate them. There's a part of this, Scott, that has to do with moral injury, and that's the idea of of betrayal.
(14:01) You know, I worked for 26 years for the CIA. I think I was involved in every covert action program in the Middle East. I did some very interesting things for the US government, always with the idea that they would have my back if I got jammed up. I just needed to get medical care when I came back and they wouldn't even do that.
(14:17) So, this moral injury, this sense of betrayal is so acute with me. Um, that's something that I can never forgive them for. >> Mark Polyopoulos rose to an executive level at the CIA, about the equivalent of a three-star general. He was awarded a top decoration for service in 2017. He says he was overwhelmed in a hotel room in Moscow. >> I woke up in the middle of the night.
(14:44) Uh it was a no I didn't hear any sound, but I woke up with incredible vertigo. Um the room was spinning. Uh I had a blinding headache. I had tonitis ringing in my ears and I felt like I was going to be physically sick. It was a terrifying feeling where I lost control. uh you know something that seriously happened to me and I remember feeling you know that this is so unusual.
(15:04) I'd been shot at in places like Iraq and Afghanistan. I'd been in physical danger but this was terrifying. He was treated for vertigo, migraines, loss of vision and trouble with memory and concentration. Disabled, he retired. Later in 2023, his own agency was among those that concluded it is very unlikely that he and the others were attacked by an adversary, which of course to me is a betrayal because CIA is supposed to be about putting people first and they did not.
(15:37) >> Are you saying this is a cover up? >> This is a massive CIA cover up and I and I'll say I say this with great regret. It's an organization that I loved. I believe in the mission. I was really good at this job. Uh to this day I want to see the CIA um operate in a strong and effective manner. >> Polyropolis and other victims have been doubted for years.
(15:59) Some in the CIA believe that a microwave weapon must be the size of a truck and so not plausible. But that changed dramatically in 2024. Three independent sources from different agencies tell us that undercover Homeland Security agents purchased a miniaturized microwave weapon from a complex Russian criminal network. It's classified.
(16:27) We didn't see it, but it has been described to us. We are told it doesn't look anything like a gun. It's designed to be concealed and small enough to be carried by a person. It is silent and doesn't create heat like a microwave oven. Our sources say the device is programmable for different scenarios and can be operated by remote control.
(16:51) The range of the beam is several hundred ft. It can penetrate windows and drywall. The vital components were made in Russia. Our sources say the key is not the hardware, but the software. The programming shapes a unique electromagnetic wave that rises and falls abruptly and pulses rapidly. >> Pulsed microwave radiation. >> Just what Dr.
(17:19) David Realman's investigations predicted. He wouldn't talk about classified information in our interview, but his research found that Russian scientists had been perfecting the concept for decades. >> And what the Russians spoke about was the importance of the energy being pulsed in order to have biological effects on humans.
(17:41) When you produce pulses like this, you can actually stimulate electrically active tissue like brain tissue and the heart for that matter, mimicking what the brain normally does, but now you're driving it with your pulses from the outside. >> An ideal stealth weapon. >> Ideal. Ideal. Because literally the person feels as if this is in my head.
(18:08) Our confidential sources tell us the still classified weapon has been tested in a US military lab for more than a year. Tests on rats and sheep show injuries consistent with those seen in humans. Also, as a separate part of the investigation, security camera videos have been collected that show Americans being hit.
(18:35) The videos are classified, but they were described to us. In one, a camera in a restaurant in Istanbul captured two FBI agents on vacation sitting at a table with their families. A man with a backpack walks in and suddenly everyone at the table grabs their head as if in pain. Our sources say another video comes from a stairwell in the US embassy in Vienna.
(19:03) The stairs lead to a secure facility. In the video, two people on the stairs suddenly collapse. Those videos and the weapon were among the reasons the Biden administration summoned about half a dozen victims to the White House with about two months left in the president's term. >> I remember the day well because I helped to organize the meeting.
(19:28) >> By that time, Dr. Realman was a White House adviser. These folks in the Biden White House believed these people and believed that their injuries were not caused by known medical or environmental conditions the way the CIA was asserting, which again to me was just egregious. Some of the specific explanations the CIA had used were just crazy.
(19:50) >> A highlevel CIA source has told us, and this is a direct quote, "This is the biggest cover up I've seen in my adult life." End quote. Do you believe it was a cover up? >> Yes, I do. Um through a variety of purposes and means, not necessarily as a um pre-planned strategic um operation, but in essence, it arrives at the same result.
(20:26) >> Help me understand what you think the motive could have been. We want this to go away so that we can resume normal operations. They had dug in opinions going back years about the plausibility of a non-therrmal microwave mechanism. In fact, when we began our work, we were briefed by their experts and told nothing in the scientific literature will support the idea that microwave energy can do things like this.
(21:00) >> They had made up their minds. >> It seems they had. And it almost seems as though consistency um was more important than objectivity. One is that >> retired CIA officer Mark Polymoropoulos was in that White House meeting. >> And so what the Biden administration was telling us is that something had changed. New intelligence had come in.
(21:27) Now I don't have a security clearance and this was an unclassified meeting. So they could not put forward that this was based on new intelligence, but it was clear to me that that's what they were insinuating. >> Dr. Paul Friedrix brought a message to the meeting. He's a retired major general, formerly one of the Pentagon's top doctors.
(21:45) >> He said very clearly, "I'm sorry. I want to apologize to you. I've never seen in 30 plus years of of practicing military medicine victims treated in such a terrible manner and I just want to offer my apologies to you." >> What did that mean to you? >> I have chills now thinking about it. I had chills then.
(22:03) It was an indication that at least some people in the Biden administration and the Biden White House believed us. Any American would be embarrassed if you were to see how these people have been treated and then to be dismissed this way as malinger or people who are manufacturing things for some other purpose. It's insulting. >> Our sources tell us that the Biden White House wrote a public statement backing the victims but never released it.
(22:34) So far, the Trump administration has not changed the words in the 2023 intelligence assessment that it is very unlikely the victims were attacked. But our sources also tell us the Trump administration has briefed top intelligence officials in Congress and shown them a classified picture of the weapon.
(22:59) We're told at the Pentagon, people who had investigated the attacks for the Department of Defense have been moved to a unit that develops new weapons. >> It was apparent to me that they did not take this issue seriously. >> Looking back, the CIA officer who quit the investigation in disgust told us in his view the CIA was careless against a ruthless adversary.
(23:26) If there was a foreign adversary, Russia in particular, would you think of this operation as a success? >> Absolutely. From an intelligence perspective, this would be a resounding success. Let's say one of these cases was real and it created all this fear, paranoia, anxiety here in the United States and overseas. The impact of that is astronomical and it's something you can't almost even calculate.
(23:54) And I don't think if it was a state actor, if it was the Russians, which I believe it is, I don't think that is something they would have put into their calculus that it would have gotten this big. And I think they saw the fear, the paranoia that it created. And I think that's why it continued to happen for that period of time over a span of a year.
(24:08) >> Across the previous stories that we've done on this subject, I have had the same question. I think you are the first person who could plausibly answer the question and that is why why would the government want to bury this? I think it comes down to a political question. I mean if we acknowledge that this was a state actor that was doing this.
(24:29) It is essentially a declaration of war against the United States which has to have a response from the United States government. In my opinion, I don't know that the appetite was there to respond to the Russians at that time. >> In our 2024 story, a collaboration with Russian dissident magazine the insider. We found evidence of Russian involvement.
(24:55) When this wife of a Justice Department official was seriously wounded overseas, an agent of Russian intelligence was in her vicinity. for this report. The Department of Defense declined to comment. The Office of the Director of National Intelligence, which oversees 18 agencies, including the CIA, told us that a new review of AHI will be comprehensive and complete, and we remain committed to delivering the truth, >> and I really hope that they do.
(25:29) The victims are waiting, including Chris and Heidi, who told us at the beginning of our story of being attacked five times in their home. >> I think it's time we as a country uh come to grips with the fact that the game has changed. Our adversaries are now able to reach out and touch us here in the United States, specifically at our homes.
(25:53) >> What do you believe the government owes you? I would say that for me and my military brothers and sisters who are hurt, uh, being issued a Purple Heart, uh, is acknowledgment of our sacrifice, uh, to the country. Um, and the sacrifice we made that affects not only us but also our families. The sources who informed our reporting told us the classified mission to obtain the microwave weapon points to a troubling reality.
(26:30) They say there are likely many of these devices. And if undercover agents could purchase one from gangsters, then the Russians have lost control of a stealth weapon that could be used by anyone anywhere. Should military victims of Havana syndrome receive the Purple Heart? >> It's not something that should be all too controversial.
(26:56) at 60 minutesovertime.com.




